% TOP_PROBLEM: {"id": 30007411, "title": "Quenched critical Ising interfaces under weak bond disorder", "category_id": 19, "category": {"id": 19, "name": "probability", "display_name": "Probability", "description": "Problems involving probability theory, stochastic processes, and random structures.", "slug": "probability", "order_index": 19, "created_at": "2026-07-31T15:26:25.671Z"}, "status": "open"}
% FIRST_POSTED: 2026-10-06
% ATTEMPT_STATUS: unresolved
% OPENAI_CLAIM: {"schema_version":1,"source_repo":"https://github.com/openai/math","source_commit":"adc7f1241b42e322a6451854ab7e4b4c146bf78a","release_date":"2026-10-06","match":"partial","resolution":"partial","independently_reviewed":false,"second_pass":"n/a","families":[{"family":"218","title":"Conformal universality for weakly interacting and random-bond Ising models","subject":"Probability and statistical mechanics","manuscripts":[{"title":"Quenched SLE ₃ Universality for the Weak Random-Bond Ising Model","pdf_path":"preprints/Quenched-SLE3-Universality-for-the-Weak-Random-Bond-Ising-Model-October-5-2026/quenched-sle3-weak-random-bond-ising.pdf","readme_path":"preprints/Quenched-SLE3-Universality-for-the-Weak-Random-Bond-Ising-Model-October-5-2026/README.md","date":"2026-10-05","theorem_ref":"thm:main (sections/introduction.tex:160)","statement_sources":[{"path":"preprints/Quenched-SLE3-Universality-for-the-Weak-Random-Bond-Ising-Model-October-5-2026/build/sections/introduction.tex","line":160,"label":"thm:main"}]},{"title":"Quenched SLE ₃ limits for general weak random-bond Ising models","pdf_path":"preprints/Quenched-SLE3-limits-for-general-weak-random-bond-Ising-models-October-5-2026/general-weak-random-bond-ising.pdf","readme_path":"preprints/Quenched-SLE3-limits-for-general-weak-random-bond-Ising-models-October-5-2026/README.md","date":"2026-10-05","theorem_ref":"thm:main (sections/introduction.tex:86)","statement_sources":[{"path":"preprints/Quenched-SLE3-limits-for-general-weak-random-bond-Ising-models-October-5-2026/build/sections/introduction.tex","line":86,"label":"thm:main"}]}],"lean":{"doc_path":"lean/docs/218.md","comparators":[],"declarations":[],"covers_main_theorem":false},"reasoning_trace":null}]}
\documentclass[11pt]{article}
\usepackage[margin=25mm]{geometry}
\usepackage{fontspec}
\setmainfont{Latin Modern Roman}
\usepackage{amsmath,amssymb,mathtools}
\usepackage{unicode-math}
\setmathfont{Latin Modern Math}
\usepackage{enumitem}
\usepackage{xurl,hyperref,bookmark}
\hypersetup{colorlinks=true,urlcolor=blue}
\setcounter{secnumdepth}{0}
\setlength{\parindent}{0pt}
\setlength{\parskip}{5pt}
\setlength{\emergencystretch}{3em}
\begin{document}
\section*{Quenched critical Ising interfaces under weak bond disorder}

\subsection{OpenAI's claim}
The linked manuscripts present the following claims by OpenAI: For small binary ferromagnetic bond disorder, critical Ising interfaces are claimed to converge quenched in probability to chordal $\mathrm{SLE}_3$; The quenched $\mathrm{SLE}_3$ conclusion is also claimed for any fixed bounded nondegenerate mean-zero iid bond law at sufficiently small disorder strength. Together they provide the following support within the scopes described below.

\subsection{Mathematical statement of the claimed result}
\paragraph{Result 1.} On square-lattice bonds let $J_e=1+\varepsilon\xi_e$, with independent fair signs $\xi_e$. Define the physical critical inverse temperature by
\[
 \beta_c(\varepsilon)=\inf\left\{\beta\ge0:
 \mathbb E_{\mathrm{env}}\lim_{m\to\infty}\langle\sigma_0\rangle^+_{[-m,m]^2,\beta,J}>0\right\},
\]
normalizing each Gibbs law before environmental averaging. Use a bounded Jordan domain with two marked boundary points and deterministic square-lattice approximations with uniformly converging marked boundary homeomorphisms. Their contour space has precisely two odd-degree dual boundary marks, completed by exterior stubs and a finite exterior contour with the same odd-degree marks. The conditional contour law is proportional to
\[
 \exp\left[-2\beta_c(\varepsilon)\sum_{e\in P}J_{e^\bullet}\right],
\]
where $e^\bullet$ crosses the dual edge $e$; bonds crossing the spin-domain boundary are included. Resolve degree four by the north-east/south-west rule and let $Q_{n,\omega}$ be the physical oriented interface law. There exists $\varepsilon_0>0$ such that for fixed $0<\varepsilon<\varepsilon_0$ and each $\eta>0$,
\[
 \mathbb P_{\mathrm{env}}\{d_{\mathrm{BL}}(Q_{n,\omega},\mathrm{SLE}_3(D;a,b))>\eta\}\to0.
\]
The metric uses uniform curve distance modulo increasing reparametrization, and the environment remains fixed while sampling each interface. The completion satisfies $P_o\subset E\setminus E_V$, where $E_V$ comprises dual edges crossing bonds incident to the spin domain $V$. The prescribed interior contour class is required to be nonempty.

\paragraph{Result 2.} Let $\rho$ be a probability law supported on $[-1,1]$ with mean zero and positive variance, take independent $\xi_e\sim\rho$, and set $J_e=1+\varepsilon\xi_e$. Define
\[
 \beta_c^\rho(\varepsilon)=\inf\left\{\beta\ge0:
 \mathbb E_\rho\lim_{m\to\infty}\langle\sigma_0\rangle^+_{[-m,m]^2,\beta,J}>0\right\}.
\]
On deterministic filled square-lattice Jordan approximations to any bounded marked Jordan domain $(D;a,b)$, use Dobrushin signs and the north-east/south-west contour resolution. Assume mesh tends to zero and marked boundary homeomorphisms tend uniformly to the identity. There is $\varepsilon_0(\rho)>0$ such that for fixed $0<\varepsilon<\min(1,\varepsilon_0(\rho))$,
\[
 d_{\mathrm{BL}}(Q_{n,\omega},\mathrm{SLE}_3(D;a,b))\xrightarrow{\mathbb P_\rho}0,
 \qquad
 \beta_c^\rho(\varepsilon)=\tfrac12\log(1+\sqrt2)+O_\rho(\varepsilon^2).
\]
Each $Q_{n,\omega}$ uses its own Gibbs normalization and uniform oriented-curve topology. The disorder bound depends on $\rho$ but is independent of domain and approximation.

\subsection{Outline of the argument}
Disorder comparison is performed for normalized Gibbs laws in a fixed environment before averaging. Protected local interface estimates and control of the disorder-induced critical shift reduce the conditional curve laws to the critical reference; tightness then identifies the full SLE law.

The binary-disorder comparison is extended by controlling general bounded bond increments. Uniform critical-reference and conditional-interface estimates identify the spontaneous-magnetization threshold and show that environment-dependent errors vanish in probability for the complete curve law.

\subsection{Scope relative to this problem}
The selected papers give the quenched interface limit first for iid fair binary bonds, then for every bounded nondegenerate mean-zero iid bond law, with the stated critical-temperature expansion. They do not cover all stationary finite-dependent environments quantified by the target.

\subsection{Formal verification}
For family 218, the formal scope is a finite-graph likelihood comparison through the zero-field FK connection probability after common pins are deleted. Stopping-band resolution and the universality companions are excluded. No Comparator is assigned to the selected manuscript here; the family page selects other papers or supporting statements. This site has not executed the Lean code. The scope descriptions do not certify the remaining manuscript arguments.

\subsection{Status caveat}
OpenAI attributes these manuscripts to work by an internal model. Its release cautions that arguments without formalization can contain mistakes. This catalogue has not independently checked the proofs; the labels report the mathematical claims attributed to OpenAI.

\subsection{Sources}
\begin{itemize}
\item OpenAI, \emph{Quenched SLE \({}_{3}\) Universality for the Weak Random-Bond Ising Model}, 2026-10-05; citation key \texttt{OAI:Quenched-SLE3-Universality-for-the-Weak-Random-Bond-Ising-Model-October-5-2026}. \href{https://github.com/openai/math/blob/adc7f1241b42e322a6451854ab7e4b4c146bf78a/preprints/Quenched-SLE3-Universality-for-the-Weak-Random-Bond-Ising-Model-October-5-2026/quenched-sle3-weak-random-bond-ising.pdf}{PDF (GitHub)}; \href{https://github.com/openai/math/raw/adc7f1241b42e322a6451854ab7e4b4c146bf78a/preprints/Quenched-SLE3-Universality-for-the-Weak-Random-Bond-Ising-Model-October-5-2026/quenched-sle3-weak-random-bond-ising.pdf}{PDF (direct)}; \href{https://github.com/openai/math/blob/adc7f1241b42e322a6451854ab7e4b4c146bf78a/preprints/Quenched-SLE3-Universality-for-the-Weak-Random-Bond-Ising-Model-October-5-2026/README.md}{Manuscript page and citation}; \href{https://github.com/openai/math/blob/main/preprints/Quenched-SLE3-Universality-for-the-Weak-Random-Bond-Ising-Model-October-5-2026/quenched-sle3-weak-random-bond-ising.pdf}{Latest version}; Statement source: \href{https://github.com/openai/math/blob/adc7f1241b42e322a6451854ab7e4b4c146bf78a/preprints/Quenched-SLE3-Universality-for-the-Weak-Random-Bond-Ising-Model-October-5-2026/build/sections/introduction.tex\#L160}{thm:main}.
\item OpenAI, \emph{Quenched SLE \({}_{3}\) limits for general weak random-bond Ising models}, 2026-10-05; citation key \texttt{OAI:Quenched-SLE3-limits-for-general-weak-random-bond-Ising-models-October-5-2026}. \href{https://github.com/openai/math/blob/adc7f1241b42e322a6451854ab7e4b4c146bf78a/preprints/Quenched-SLE3-limits-for-general-weak-random-bond-Ising-models-October-5-2026/general-weak-random-bond-ising.pdf}{PDF (GitHub)}; \href{https://github.com/openai/math/raw/adc7f1241b42e322a6451854ab7e4b4c146bf78a/preprints/Quenched-SLE3-limits-for-general-weak-random-bond-Ising-models-October-5-2026/general-weak-random-bond-ising.pdf}{PDF (direct)}; \href{https://github.com/openai/math/blob/adc7f1241b42e322a6451854ab7e4b4c146bf78a/preprints/Quenched-SLE3-limits-for-general-weak-random-bond-Ising-models-October-5-2026/README.md}{Manuscript page and citation}; \href{https://github.com/openai/math/blob/main/preprints/Quenched-SLE3-limits-for-general-weak-random-bond-Ising-models-October-5-2026/general-weak-random-bond-ising.pdf}{Latest version}; Statement source: \href{https://github.com/openai/math/blob/adc7f1241b42e322a6451854ab7e4b4c146bf78a/preprints/Quenched-SLE3-limits-for-general-weak-random-bond-Ising-models-October-5-2026/build/sections/introduction.tex\#L86}{thm:main}.
\item Family 218: \href{https://github.com/openai/math/blob/adc7f1241b42e322a6451854ab7e4b4c146bf78a/lean/docs/218.md}{Lean scope}.
\item \href{https://github.com/openai/math/blob/adc7f1241b42e322a6451854ab7e4b4c146bf78a/CONTENTS.md}{OpenAI catalogue entry}.
\end{itemize}
\end{document}
